\documentclass[10pt,a4paper]{amsart}
\usepackage[margin=19mm]{geometry}
\usepackage{amsmath,amssymb,mathtools}
\usepackage[scr=rsfs,cal=euler]{mathalfa}
\usepackage{xfrac}
\usepackage[unicode,hidelinks,bookmarks=false]{hyperref}
\newcommand{\ie}{i.e.\ }
\newcommand{\cf}{cf.\ }
\newcommand{\resp}{respectively\ }
\newcommand{\Subsection}{\subsection}
\providecommand{\nobreakdash}{\nobreak\hbox{-}}
\setlength{\emergencystretch}{2em}
\multlinegap0pt
\providecommand{\texorpdfstring}[2]{#1}
\newcommand{\st}{\textrm{st}}
\newcommand{\lk}{\textrm{lk}}
\newcommand{\cat}{\textrm{CAT}}
\newcommand{\dime}{\textrm{dim}}
\newcommand{\degg}{\textrm{deg}}
\renewcommand{\sigma}{g}
\renewcommand{\varphi}{h}
\newcommand{\qq}{\mathbb{Q}}
\newcommand{\x}{\mathcal{X}}
\newcommand{\kp}{\mathcal{K}}
\newcommand{\kun}{\mathfrak{K}}
\newcommand{\N}{\mathbb{N}}
\newcommand{\R}{\mathbb{R}}
\newcommand{\Z}{\mathbb{Z}}
\newcommand{\cc}{\mathrm{c}}
\newcommand{\aut}{\textrm{Aut}}
\newcommand\restr{\raisebox{-0.3ex}{$|$}\raisebox{0.3ex}{}}
\newtheorem{cdrthm}{Theorem}[section]
\theoremstyle{plain}
\newtheorem{teo}[cdrthm]{Theorem}
\newtheorem{lema}[cdrthm]{Lemma}
\newtheorem{prop}[cdrthm]{Proposition}
\newtheorem{coro}[cdrthm]{Corollary}
\theoremstyle{definition}
\newtheorem{defi}[cdrthm]{Definition}
\newtheorem{ej}[cdrthm]{Example}
\theoremstyle{remark}
\newtheorem{obs}[cdrthm]{Remark}
\title[Smallest posets with given cyclic automorphism group]{Exact minimum size of finite posets with a specified finite cyclic automorphism group: original Corollary 4.2}
\author{Jonathan Ariel Barmak and Agust\'in Nicol\'as Barreto}
\date{Source: Algebraic Combinatorics 7 (2024), 1307--1318; DOI 10.5802/alco.372; publisher version}
\begin{document}
\maketitle
\noindent\textit{Editorial scope.} The counted unit is original Corollary 4.2 with the author's complete proof, preceded here by the complete source-local core it uses: the poset and cycle terminology, the extremal constructions giving the upper bound, the structural lemmas on cycles of a generator, and the prime-power weight theorem with its complete proof. The exact-divisibility hypothesis and the exceptional minus-one constant are unchanged, every retained passage is native author text, and printed slips are kept literal. The five original illustrations are not redistributed; their captions, labels and cross-references are retained.
\noindent\textit{Original introduction passage defining $\beta$ (uncounted context).}
In parallel, the analogous problem was studied for partially ordered sets. In 1946 G.~Birkhoff~\cite{Bir} proved that for any finite group $G$ there is a poset of $|G|(|G|+1)$ points and automorphism group isomorphic to $G$. Then Frucht~\cite{Fru50} improved this to $(d+2)|G|$ points. In 1980 Babai~\cite{Bab2} proved that $3|G|$ points are enough. However, the smallest number $\beta (G)$ of points of a poset with an arbitrary finite abelian group $G$ of automorphisms has not yet been determined. In this paper we compute $\beta (G)$ for every finite cyclic group $G$. 

\setcounter{section}{1}
\section{Construction of the examples} \label{sectionexamples}

A poset is a set with a partial order $\le$. The elements of the underlying set of a poset are called points. All posets are assumed to be finite, that is, their underlying set is finite. A subposet of a poset $P$ is a subset of the underlying set with the inherited order. A chain in a poset is a subset of pairwise comparable points. The height of a poset is the maximum cardinality of its chains minus one. A poset, or a subset of a poset, is said to be discrete if no two points are comparable. If $P$ is a poset and $x,y\in P$, we write $x<y$ if $x\le y$ and $x\neq y$. We say that $y$ covers $x$ if $x<y$ and there is no $x<z<y$. The edges of $P$ are the pairs $(x,y)$ such that $y$ covers $x$. The Hasse diagram of $P$ is the digraph whose vertices are the points of $P$ and the edges are the edges of $P$. If the orientation of an arrow is not indicated in the graphical representation of the Hasse diagram, we assume it points upwards. A morphism $P\to Q$ of posets is an order-preserving map, \ie a function $f$ between the underlying sets such that for every pair $x,y\in P$ with $x\le y$ we have $f(x)\le f(y)$. If $P$ is a poset, since it is finite, an automorphism of $P$ is just a permutation of the underlying set which is a morphism. Given an automorphism $g$ of a poset $P$, we say that a subset $A$ of the underlying set of $P$ is invariant or $g$-invariant if $g(A)=A$. In this case, $g$ induces an automorphism on the subposet with underlying set $A$.  

\begin{defi}
Define $b(1)=0$, $b(2)=1$, $b(3)=b(4)=b(5)=b(7)=3$. For any other prime power $p^r$, define $b(p^r)=2$. 
\end{defi}

\begin{defi}
Let $P$ be a poset of height 1 and let $A$ be the set of minimal points and~$A'$ the set of maximal points. Given $a\in A$ we define its \emph{double neighborhood} $B(a)$ as the set of those $b \in A$ such that $\# (P_{> a}\cap P_{> b}) \ge 2$, that is, there are at least two points in $A'$ greater than both, $a$ and $b$. The \emph{reduced double neighborhood} of $a\in A$ is~$\hat{B}(a)=B(a)\smallsetminus \{a\}$. Given $k\ge 1$, we say that two points $a,b\in A$ are \emph{$k$-adjacent} if~$\#(B(a)\cap B(b))=k$, and they are \emph{reduced $k$-adjacent} if~$\#(\hat{B}(a)\cap \hat{B}(b))=k$.
\end{defi}

We denote by $\Z_n=\{0,1,\ldots, n-1\}$ the additive group of integers modulo $n$.

\begin{prop} \label{ejemplos}
Let $n=p^r$, where $p\ge 2$ is a prime and $r\ge 0$. Then there exists a poset $P$ with $b(n)n$ points and automorphism group $\aut(P)$ isomorphic to $\Z_{n}$.  
\end{prop}

\begin{proof}
For $n=1$ we take the empty poset and for $n=2$ we take the discrete poset on $2$ points. If $n=3,4,5,7$ we use the following well-known general construction~\cite{Fru50}: $P=\Z_{n}\times \{0,1,2\}$ with the order $(i,2)>(i,1)>(i,0)<(i+1,2)$ for every $i\in \Z_{n}$. It is easy to see that such poset satisfies $\aut(P) \simeq \Z_{n}$. Suppose then that $n\ge 8$. We take two copies of $\Z_{n}$: $A=\Z_{n}=\{0,1,\ldots, n-1\}$ and $A'=\{0',1',\ldots, (n-1)'\}$. Let $S=\{0,1,2,4\}\subseteq \Z_{n}$. For $i\in A$ and $j'\in A'$ we set $i<j'$ if $j-i\in S$. Any two elements in the same copy of $\Z_n$ are not comparable (see Figure~\ref{fig0124}). We will prove that the automorphism group of this poset $P$ is $\Z_{n}$. It is clear that $G=\Z_n$ acts regularly on each copy of $\Z_n$ by addition, and this gives a faithful action $G\to \aut(P)$ on $P$. So~$G$ can be seen as a subgroup of $\aut(P)$. Since each automorphism of $P$ maps $0\in A$ to another minimal element of $P$, the order of the $\aut(P)$-orbit of $0\in P$ is $n$. If we prove that the $\aut(P)$-stabilizer of $0\in P$ is trivial, then $|\aut(P)|=n$, so $\aut(P)$ is isomorphic to $G$. Let $\varphi \in \aut(P)$ be such that $\varphi (0)=0$.

\begin{figure}[h] 
\begin{center}


\caption{The Hasse diagram of $P$ for $n=8$.}\label{fig0124}
\end{center}
\end{figure}

Since $\varphi$ is an automorphism, for every $i\in A$ one has $B(\varphi(i))=\varphi (B(i))$ and $\hat{B}(\varphi(i))=\varphi (\hat{B}(i))$. Thus, $\varphi$ preserves $k$-adjacency and reduced $k$-adjacency. Suppose first that $n\ge 9$. Then for each $i\in A$, $B(i)=\{i-2,i-1,i,i+1,i+2\}$. It is easy to see that $i,j$ are $4$-adjacent if and only if $i-j=\pm 1$. Thus, $\varphi$ induces an automorphism of the cyclic graph on $A$ with edges given by $4$-adjacency. Since $\varphi(0)=0$, $\varphi$ is either the identity $1_{\Z_n}$ or $-1_{\Z_n}$ (\ie the map $i\mapsto -i$). The second case cannot occur as $\{0,2,3,4\}$ has an upper bound while $\{0,-2,-3,-4\}$ does not. Thus every point of $A$ is fixed by $\varphi$. If $j'\in A'$, then $j'$ is the unique upper bound of $\{j,j-1,j-2,j-4\}$. Thus $\varphi(j')=j'$. This proves that $\varphi=1_P$. 

Finally, suppose $n=8$. Given $i\in A$, we have now $\hat{B}(i)=\{i-2,i-1,i+1,i+2, i+4\}$ and $i,j\in A$ are reduced $4$-adjacent if and only if $i-j=\pm 3$. Thus, $\varphi$ induces an automorphism in the cyclic graph on $A$ with edges given by reduced $4$-adjacency. Then $\varphi=1_{\Z_n}$ or $-1_{\Z_n}$. The second case cannot occur for the same reason as before. Since each point in $A'$ is determined by the set of smaller points, $\varphi=1_P$.    
\end{proof}

\begin{ej} \label{ejemplo12}
There exists a poset $P$ with $20$ points and automorphism group isomorphic to $\Z_{12}$.

Take two copies $A=\{0,1,2,3,4,5\}$, $A'=\{0',1',2',3',4',5'\}$ of $\Z_6$ and two copies $B=\{0'',1'',2'',3''\}$, $B'=\{0''',1''',2''',3'''\}$ of $\Z_4$. The underlying set of $P$ is the union of these four sets. Let $S=\{0,1,3\}\subseteq \Z_6$, $T=\{0,1\}\subseteq \Z_4$. Define the following order in $P$: $i<j'$ if $j-i\in S$, $i''<j'''$ if $j-i\in T$, $i'''<j'$ if $j-i$ is even, $i''<j$ if $j-i$ is even, $i''<j'$ for every $i,j$ (see Figure~\ref{figveinte}). 

\begin{figure}
\begin{center}


\caption{A poset $P$ of $20$ points and $\aut(P)\simeq \Z_{12}$.}\label{figveinte}
\end{center}
\end{figure}

It is clear that $G=\Z_{12}$ acts in each copy of $\Z_6$ and of $\Z_4$ by addition. This induces a faithful action of $G$ on $P$. If $\varphi\in \aut(P)$, $\varphi(0'')$ must be a minimal point $i''$ and $\varphi(0')$ must be a maximal point $j'$. However $i,j$ cannot have different parity. Indeed, among the points $0,2,4,0''',1'''$ which cover $0''$, there are just two $0,0'''$ smaller than~$0'$. However, if $i\in \Z_4$ and $j\in \Z_6$ have different parity, among the points covering $i''$ ($k\in A$ with $k\equiv i (2)$ and $i''', (i+1)'''$) there are three smaller than $j'$: both $j-1, j-3$, and one of $i''', (i+1)'''$. Thus $i\equiv j(2)$, which implies that the $\aut(P)$-orbit of the set $\{0',0''\}$ has at most $12$ elements. If we prove that the $\aut(P)$-stabilizer of $\{0',0''\}$ is trivial, then $|\aut(P)|\le 12=|G|$, so $\aut(P)$ is isomorphic to $G$.
Let $\varphi$ be an automorphism of $P$ which fixes $0'$ and $0''$.

Note that $2''$ is the unique minimal point different from $0''$ which is covered by three points that cover $0''$. Thus $\varphi(2'')=2''$. Now, the points of $B'$ are the unique points of $P$ which cover exactly one of $0'',2''$. Thus $B'$ is invariant. This implies that $\varphi$ restricts to an automorphism of the subposet $R$ with underlying set $B\cup B'$ and of the subposet $Q$ with set $A\cup A'$.
Since $R$ is a cycle, there are only two automorphisms of $R$ fixing $0''$. One is the identity and the other maps $0'''$ to $1'''$. However, $0'''<0'$ while $1'''\nless 0'$. Thus $0'''$ is fixed by $\varphi$ and then $\varphi$ is the identity of $R$. 

Suppose that $i'\in A'$ is a fixed point. Among the points $i,i-1,i-3$ in $A$ covered by $i'$, only $i-1$ and $i-3$ share a lower bound. Thus $\varphi(i)=i$. Similarly, among the points $(i-4)',(i-2)',(i-1)'$ of $A'$ not covering $i$, only $(i-4)'$ and $(i-2)'$ share a lower bound in $B'$. Thus $(i-1)'$ is fixed. In conclusion, we showed that $i'$ fixed implies that both $i$ and $(i-1)'$ are fixed. Since $0'$ is fixed, this implies that every point of $A$ and of $A'$ is fixed. Thus $\varphi=1_P$.  

\end{ej}

We say that a prime power $p^r$ ($r\ge 1$) exactly divides an integer $n$, and write $p^r\parallel n$, if $p^r| n$ and $p^{r+1}\nmid n$.

\begin{teo} \label{teoejemplos}
Let $n=p_1^{r_1}p_2^{r_2}\ldots p_k^{r_k}$ where the $p_i$ are pairwise different primes and~$r_i\ge 1$ for every $i$. Then there exists a poset with automorphism group isomorphic to $\Z_n$ and $\sum\limits_{i=1}^k b(p_i^{r_i})p_i^{r_i}-1$ points if $3\parallel n$ and $4\parallel n$, and with $\sum\limits_{i=1}^k b(p_i^{r_i})p_i^{r_i}$ points otherwise.
\end{teo}

\begin{proof}
By Proposition~\ref{ejemplos}, for each $1\le i\le k$ there exists a poset $P_i$ with $b(p_i^{r_i})p_i^{r_i}$ points and $\aut(P_i)\simeq \Z_{p_i^{r_i}}$. The non-Hausdorff join or ordinal sum $P=P_1\oplus P_2\oplus \ldots \oplus P_k$ is constructed by taking a copy of each poset and keeping the given ordering in each copy, while setting $x<y$ for each $x\in P_i$ and $y\in P_j$ if $i<j$. Since each automorphism of $P$ preserves heights (the maximum length of a chain with a given maximum element), it restricts to automorphisms of each $P_i$. Thus $\aut(P)=\aut(P_1)\oplus \aut(P_2) \oplus \ldots \oplus \aut(P_k)=\Z_n$. If $p_i^{r_i}=3$ and $p_j^{r_j}=4$, instead of $P_i$ and $P_j$ we take the poset in Example~\ref{ejemplo12} of $20=b(3)3+b(4)4-1$ points and automorphism group $\Z_{12}$.
\end{proof}

\section{Lemmas} \label{sectionlemmas}


Let $X$ be a finite set, $n\ge 1$ and $x_0,x_1,\ldots, x_{n-1}$ pairwise different elements of $X$. The cycle $\alpha =(x_0,x_1,\ldots , x_{n-1})$ is the permutation which maps $x_i$ to $x_{i+1}$ (indices considered modulo $n$) and fixes every other point of $X$. The number $n$ is the order or length of the cycle, which we denote by $|\alpha|$. A cycle of order $n$ is also called an $n$-cycle. A cycle $\alpha$ is non-trivial if $|\alpha|\ge 2$. The representation $(x_0,x_1,\ldots, x_{n-1})$ of a non-trivial $n$-cycle is unique up to cyclic permutation of the $n$-tuple $x_0,x_1,\ldots,x_{n-1}$. The underlying set of a non-trivial cycle $(x_0,x_1,\ldots , x_{n-1})$ is $\{x_0,x_1,\ldots , x_{n-1}\}$. Many times we will identify a non-trivial cycle with its underlying set. Two non-trivial cycles are disjoint if their underlying sets are. Any permutation $g$ of $X$ can be written as a composition $\alpha_1 \alpha_2 \ldots \alpha_k$ of pairwise disjoint non-trivial cycles. This representation is unique up to reordering of the cycles.
If a cycle $\alpha$ appears in the factorization of $g$, we say that $\alpha$ is contained in $g$ and write $\alpha \in g$. The orbits of $g$, or of the action of the cyclic group $\langle g \rangle$ on $X$, are the underlying sets of the cycles in $g$ and the singletons consisting of fixed points. Disjoint non-trivial cycles commute. Thus, if $g$ is a composition $\alpha_1 \alpha_2 \ldots \alpha_k$ of pairwise disjoint non-trivial cycles and $m\in \Z$, then $g^m=\alpha_1^m \alpha_2^m \ldots \alpha_k^m$. If $\alpha$ is a cycle of length $n$ and $m\in \Z$, the permutation $\alpha^m$ is a composition of $(n,m)=$gcd$\{n,m\}$ cycles of length $\frac{n}{(n,m)}$. In particular, $\alpha^m$ is a cycle with the same underlying set as $\alpha$ if $n$ and $m$ are coprime. Moreover, the order of $g$ is the least common multiple of the lengths of its cycles and if a cycle of $g$ has order~$n$, and $m\in \Z$, then $g^m$ fixes every point of the cycle if $n|m$, and fixes no point of the cycle otherwise. 


If $g$ is an automorphism of a poset $P$, then each orbit of $g$ is discrete, as $a<b$ would imply that $a<g^k(a)$ for some $k\in \Z$ and then $\{g^{nk}(a)\}_{n\ge 0}$ would be an infinite chain. If $A$ and $B$ are two different orbits of $g$ we cannot have an element $a\in A$ smaller than another $b\in B$ and at the same time an element $b'\in B$ smaller than another $a'\in A$, as this would imply that $a<b=g^k(b')<g^k(a')$ for some $k\in \Z$, contradicting the fact that $A$ is discrete, or the antisymmetry of the order.


\begin{obs} \label{extension}
Let $P$ be a poset and let $g$ be an automorphism of $P$. Let $Q$ be the subposet of points which are not fixed by $g$. Let $A_0,A_1,\ldots, A_k$ be the orbits of the automorphism induced by $g$ on $Q$. If $h$ is an automorphism of $Q$ such that $h(A_i)=A_i$ for every $i$, then it extends to an automorphism of $P$ which fixes every element not in $Q$.

Indeed, if $x\in P\smallsetminus Q$, $y\in A_i$ and $x<y$, then $h(y)\in A_i$, so there exists $r\ge 0$ such that $g^r(y)=h(y)$. Then $x=g^r(x)<g^r(y)=h(y)$. Similarly, if $x>y$, then $x>h(y)$.
\end{obs}

\begin{lema} \label{dos}
Let $n\ge 1$ and let $p^r\neq 2$ be a prime power which exactly divides $n$. Let~$P$ be a poset with $\aut(P)$ cyclic of order $n$, and let $g$ be a generator of $\aut (P)$. Then $g$ contains at least two cycles of length divisible by $p^r$. 
\end{lema}
\begin{proof}
Since $g$ has order $n$, it contains at least one cycle $\alpha$ of length divisible by $p^r$. Assume there is no other cycle of length divisible by $p^r$. The automorphism $g^{\frac{n}{p}}$ fixes then every point not in $\alpha$. Let $x$ be an element of $\alpha$ and let $\tau$ be the transposition of the underlying set of $\alpha$ which permutes $x$ and $g^{\frac{n}{p}}(x)\neq x$. By Remark~\ref{extension}, $\tau$ extends to an automorphism $h$ of $P$ which is a transposition. But any power of $g$ either fixes each point in $\alpha$ or fixes no point of $\alpha$. Since the order of $\alpha$ is at least $p^r>2$, $h \notin \langle g\rangle=\aut(P)$, a contradiction.  
\end{proof}

If a group $G$ acts on a poset $P$, an automorphism of $P$ is said to be induced by the action if it is in the image of the homomorphism $G\to \aut(P)$. 

Let $P$ be a poset of height 1. Let $A$ be the set of minimal points and let $B$ be the set of maximal points. Suppose that for every $a\in A$ there exist $b,b'\in B$ such that $a<b$ and $a\nless b'$, and dually for every $b\in B$ there exist $a,a'\in A$ with $a<b$ and $a'\nless b$. We define the \emph{complement} $P^c$ of $P$ to be the height 1 poset with the same sets of minimal and maximal points and where $a<b$ for $a\in A, b\in B$ if and only if $a\nless b$ in $P$. Note that in this case $P$ and $P^c$ have the same automorphisms.

\begin{lema} \label{nuevofacil}
Let $p=3,5$ or $7$. Let $P$ be a poset on which $\Z_p$ acts with exactly two orbits, both of order $p$. Then there exists an automorphism of $P$ not induced by the action for which each orbit of the action is invariant.
\end{lema}

\begin{proof}
Let $g=\alpha \beta \in \aut(P)$ be the automorphism induced by a generator of $\Z_p$, where $\alpha=(0,1,\ldots, p-1)$ and $\beta=(0',1',\ldots, (p-1)')$. If no element of $\alpha$ is comparable with an element of $\beta$, then the transposition $(0,1)$ is an automorphism which is different to $g^k$ for any $k\in \Z$, that is, not induced by the action.

Without loss of generality we can assume then that $0$ and $0'$ are comparable, and moreover, that $0<0'$. Then no element in $\beta$ can be smaller than another in $\alpha$. Since~$g$ is an automorphism, $i<i'$ for every $0\le i\le p-1$. If no other pair of elements are comparable, then $(0,1)(0',1')$ is an automorphism not induced by the action (it has order $2$, for example). If $i<j'$ for every $0\le i,j\le p-1$, then $(0,1)$ satisfies the desired property. This completes the proof of the case $p=3$ by the following argument. The case we did not analyze is when $P$ has exactly $6$ edges. In that case the complement $P^c$ has only $3$ edges, so there is an automorphism of $P^c$ not induced by the action, and this is the required automorphism of $P$.


For $p=5$ we need to consider the case that $P$ has $10$ edges. By the complement argument, this will complete the $p=5$ case. So, suppose $0<k'$ for some $1\le k\le 4$ (and then $i<(i+k)'$ for every $i$, where $i+k$ is considered modulo $5$). Note that~$g^k$ is induced by another generator of $\Z_p$ and it maps $i'$ to $(i+k)'$. Thus, for each $0\le i\le 4$, $i<i'$ and $i<g^k(i')$. Therefore we can assume that $k=1$. We have then the ``symmetry about the axis $03'$'', which maps $i$ to $-i$ and $j'$ to $(1-j)'$ (see Figure~\ref{figcinco}). This is an automorphism of $P$ which is different to any power of $g$ (it has order $2$).

\begin{figure}[h] 
\begin{center}


\caption{The underlying undirected graph of a poset with $10$ points and edges $i'>i<(i+1)'$, and the axis $03'$.}\label{figcinco}
\end{center}
\end{figure}

For $p=7$, if $P$ has $14$ edges, then by the argument above we can assume $i'>i<(i+1)'$ for every $0\le i\le 6$ and there is then a symmetry about $04'$. By the complement argument it only remains to analyze the case that $P$ has exactly $21$ edges. Here $i<i',(i+k)', (i+l)'$ for certain $1\le k\neq l\le 6$ and again we can assume $k=1$ by replacing $g$ by $g^k$. Finally, by replacing $g$ by $g^{-1}$, it suffices to consider the cases $l=2,3$ and $4$ (Figure~\ref{figsietetres}).

\begin{figure}[h] 
\begin{center}



\caption{Posets with two $\Z_7$-regular orbits and $S=\{0,1,l\}$ for $l=2,3,4$.}\label{figsietetres}
\end{center}
\end{figure}

For $l=2$ we have the involution that maps $i$ to $-i$ and $j'$ to $(2-j)'$. For $l=3$ we have the following automorphism of order $3$: $(142)(356)(0'3'1')(2'4'5')$ (see Figure~\ref{triangulo}). For $l=4$, there is again the symmetry about $04'$. 
\begin{figure}[h] 
\begin{center}



\caption{The underlying graph of the poset $P$ of $14$ points and edges $i<i',(i+1)',(i+3)'$. An automorphism of order $3$ is given by a rotation of angle $\frac{2\pi}{3}$.}\label{triangulo}
\end{center}
\end{figure}
\end{proof}

\begin{lema} \label{nuevodificil}
Let $P$ be a poset on which $\Z_4$ acts with exactly two orbits of order $4$ or exactly three orbits: two  of order $4$ and one of order $2$. Then there exists an automorphism of $P$ not induced by the action for which each orbit of the action is invariant.
\end{lema}

\begin{proof}
Let $g$ be an automorphism induced by a generator of the action and suppose first that $g=(0,1,2,3)(0',1',2',3')$. If $P$ is discrete or has $16$ edges, $(0,1)$ satisfies the required conditions. If $P$ has exactly $4$ edges, then as in the proof of Lemma~\ref{nuevofacil} we can assume $i<i'$ for every $0\le i\le 3$, and $(0,1)(0',1')$ works. By the complement argument we can assume $P$ has exactly $8$ edges and that it is determined by the relations $i'>i<(i+k)'$ for some $1\le k\le 3$. The case $k=3$ reduces to the case $k=1$ by replacing $g$ by $g^3$. If $k=1$, the symmetry $(1,3)(0',1')(2',3')$ about $02$ satisfies the required conditions. If $k=2$, then $(0,2)$ works.

Suppose then that $g=\alpha \beta \gamma$ with $\alpha=(0,1,2,3)$, $\beta=(0',1',2',3')$, $\gamma=(0'',1'')$. Let $Q$ be the subposet of points in $\alpha$ and $\beta$. Since $g^2=(0,2)(1,3)(0',2')(1',3')$, every automorphism of the poset $Q$ which has $\{0,2\},\{1,3\}, \{0',2'\}, \{1',3'\}$ as invariant sets, extends to $P$ by Remark~\ref{extension}. If $Q$ is discrete or if $Q$ has $16$ edges, then $(0,2)$ is an automorphism of $Q$ which extends to $P$ and this extension is not induced by the action. If $Q$ has exactly $4$ edges, we may assume $i<i'$ for every $i$ and then $(0,2)(0',2')$ extends to an automorphism of $P$ different to any power of $g$. If $Q$ has exactly 12 edges, the complement argument can be used. Suppose then $Q$ has exactly $8$ edges. By relabelling we can assume the relations are (a) $i<j'$ for $i\equiv j (2)$ or (b) $i'>i<(i+1)'$ for every $i$. In case (a), $(0,2)$ is again an automorphism which has every nontrivial orbit of $g^2$ as an invariant set. In the rest of the proof we assume we are in case (b). 

If the points of $\gamma$ are not comparable with any point of $Q$, then the symmetry about~$02$ which maps $i$ to $-i$ and $j'$ to $(1-j)'$, is an automorphism of $Q$ which extends to~$P$, and this extension satisfies the required conditions.

By considering the opposite order, we can assume a point of $\gamma$ is comparable with a point of $\alpha$. Moreover, by relabelling if needed we can assume $0''$ is comparable with~$0$. Suppose first that $0''<0$. Since $g$ is an automorphism, then $0''<2$ and $1''<1,3$. If~$0''\nless 1$, then $0''\nless 3$ and $1''\nless 0,2$. If $0''<1$, then $0''<3$ and $1''<0,2$. In either case, the symmetry of $Q$ about $02$ extends by the identity to an automorphim of $P$ which is not induced by the action, even though this automorphism of $Q$ does not have the orbits of $g^2$ as invariant sets. Finally suppose $0''>0$. Then $0''>2$ and $1''>1,3$. We can assume no element in $\beta$ is smaller than an element in $\gamma$, by the previous case and the duality argument. Also, we cannot have an element of $\gamma$ being smaller than another $j'$ of $\beta$, since this would imply that $i<j'>i+2$, modulo 4, for certain $0\le i\le 3$, which is absurd. In any case, if $0''\ngtr 1$ or if $0''>1$, we have that the symmetry of $Q$ about $02$ extends to an automorphism of $P$. 
\end{proof}

\begin{lema} \label{lema357}
Let $p=3,5$ or $7$. Let $P$ be a poset with cyclic automorphism group of order $n\ge 1$, and let $g\in \aut(P)$ be a generator. Suppose $g$ contains a $p$-cycle $\alpha$ and a $pk$-cycle $\beta\neq \alpha$ for some $p\nmid k\ge 1$. Then it contains a third cycle whose length is divisible by $p$.
\end{lema}

\begin{proof}
Suppose $\beta=(0,1,\ldots, pk-1)$. Let $Q$ be the subposet of $P$ whose points are those of $\alpha$ and $\beta$. Assume that there is no other cycle in $g$ whose length is divisible by~$p$. In particular $p\parallel n$. Since the order of any cycle of $g$ different from $\alpha$ and $\beta$ divides $\frac{n}{p}$, the automorphism $g^{\frac{n}{p}}$ fixes every point not in $Q$. Moreover $g^{\frac{n}{p}}$ has $k+1$ orbits of order $p$, which are the underlying set of $\alpha$ and $A_i=\{0\le j\le pk-1 | \ j\equiv i (k)\}$ for $0\le i\le k-1$. In particular, by Remark~\ref{extension} every automorphism of $Q$ for which these sets are invariant extends to an automorphism of $P$.

Let $Q'$ be the subposet of $Q$ whose points are those of $\alpha$ and $A_0$. Since $g^{k}$ induces an automorphism of $Q'$ with two orbits of order $p$, by Lemma~\ref{nuevofacil} there is an automorphism $\varphi$ of $Q'$ not induced by a power of $g^{k}$ for which the underlying set of $\alpha$ and $A_0$ are invariant. We extend $\varphi$ to an automorphism $\overline{\varphi}$ of $Q$ as follows. Let $j$ be a point of $\beta$, $0\le j\le kp-1$. Let $0\le i\le k-1$ be such that $j\in A_i$. Since $p\nmid k$, there exists a unique  $0\le t\le k-1$ such that $k|j+tp$, in other words $j+tp$, considered modulo $kp$, lies in $A_0$. Then $\varphi(j+tp)\in A_0$. Define $\overline{\varphi}(j)=\varphi(j+tp)-tp\in A_i$. We claim that $\overline{\varphi}$ is an automorphism of $Q$. It is clearly bijective. Two different points of $\beta$ cannot be comparable as they are in the same orbit. Suppose $j$ in $\beta$ and $a$ in $\alpha$ are comparable, say $a<j$. Let $0\le t\le k-1$ be such that $k|j+tp$. Then $a=g^{tp}(a)<g^{tp}(j)=j+tp$. Since $\varphi$ is a morphism, $\varphi(a)<\varphi(j+tp)$. Thus $\overline{\varphi}(a)=\varphi (a)=g^{-tp}(\varphi(a))<g^{-tp} (\varphi(j+tp))=\varphi(j+tp)-tp=\overline{\varphi}(j)$. Since the underlying set of $\alpha$ and each $A_i$ are $\overline{\varphi}$-invariant, $\overline{\varphi}$ extends to an automorphism of $P$, which must be a power $g^r$ of $g$. Since $g^r$ leaves $A_0$ invariant, in particular $r=g^r(0)\in A_0$, so $k|r$ and $\varphi$ is then induced by a power of $g^k$, a contradiction.   
\end{proof}

\begin{lema} \label{lema4}
Let $P$ be a poset with cyclic automorphism group of order $n\ge 1$, and let $g\in \aut (P)$ be a generator. Suppose that $g$ contains two $4$-cycles $\alpha, \beta$. Then it contains a third cycle of length divisible by $4$ or two more cycles of even length. 
\end{lema}

\begin{proof}
The proof is very similar to that of Lemma~\ref{lema357}, so we omit details. If $\alpha$ and~$\beta$ are the unique two cycles of even length in $g$, then by Lemma~\ref{nuevodificil} there is an automorphism~$h$ of the poset of points of these two cycles which is not induced by a power of $g$, and moreover has the underlying sets of $\alpha$ and $\beta$ as invariant sets. Since the non-trivial orbits of $g^{\frac{n}{4}}\in \aut(P)$ are the underlying sets of $\alpha$ and $\beta$, $h$ extends to an automorphism of $P$, a contradiction.  

Suppose then there exists a third cycle $\gamma=(1,2,\ldots, 2k)$ in $g$ with $k$ odd, and that there is no other cycle of even length. We define $Q$ to be the subposet whose points are those of $\alpha, \beta$ and $\gamma$. Then $g^{\frac{n}{4}}$ fixes every point not in $Q$. The other orbits of $g^{\frac{n}{4}}$ are the underlying sets of $\alpha$ and $\beta$, and $A_i=\{i,k+i\}$ for $0\le i\le k-1$. Let $Q'$ be the subposet whose points are those of $\alpha, \beta$ and $A_0$. Then $g^k$ induces an automorphism of $Q'$ and by Lemma~\ref{nuevodificil} there is an automorphism $\varphi$ of $Q'$ which is not induced by a power of $g^k$, and for which the underlying sets of $\alpha, \beta$ and $A_0$ are invariant. We extend it to an automorphism $\overline{\varphi}$ of $Q$ by defining $\overline{\varphi}(j)=\varphi(j+4t)-4t$, where $t$ is such that $k|j+4t$. Then $\overline{\varphi}$ is bijective, it is a morphism and leaves each $A_i$ invariant. It extends to an automorphism of $P$, say $g^r$. Since $g^r$ leaves $A_0$ invariant, then $k|r$, which implies that $h$ is induced by a power of $g^k$, a contradiction.
\end{proof}

\section{Weights and the lower bound}  \label{sectionweights}

Let $\sigma$ be a permutation of order $n$ of a finite set $X$. Let $\alpha$ be a cycle in $\sigma$ of length $l=p_1^{r_1}p_2^{r_2}\ldots p_k^{r_k}$, where the $p_i$ are pairwise distinct prime integers, $r_i\ge 1$ for every~$i$. For each prime power $p^r$ we will define a weight $w_{p^r}(\alpha)\in \R_{\ge 0}$ which depends on $p^r,l$ and $n$, in such a way that $\sum\limits_{p^r} w_{p^r}(\alpha)p^r= l$, where the sum is taken over all prime powers dividing $n$. In particular $\# X\ge \sum\limits_{p^r\parallel n} (\sum\limits_{\alpha \in \sigma} w_{p^r}(\alpha))p^r$. For each $l\ge 2$ we will assign the weight of every prime power $p^r$ in a cycle $\alpha$ of length $|\alpha|=l$ according to a series of rules. In every case, if the weight $w_{p^r}(\alpha)$ is not explicitly defined for some prime power, we assume it is $0$.  

\noindent \textbf{Exception 6}. Suppose $l=6$. If $3\parallel n$ then $w_3(\alpha)=2$. If $3\nparallel n$ and $2\parallel n$, then $w_{2}(\alpha)=3$. If $3\nparallel n$ and $2\nparallel n$, then $w_4(\alpha)=\frac{3}{2}$.

\noindent \textbf{Exception 12}. Suppose $l=12$. If $3\parallel n$ then $w_3(\alpha)=4$. If $3\nparallel n$, then $w_{4}(\alpha)=3$.

\noindent \textbf{Exception 10-14}. Suppose $l=2p$ for $p=5$ or $7$. If $2\parallel n$, $w_2(\alpha)=1$. Otherwise $w_4(\alpha)=\frac{1}{2}$. In any case $w_p(\alpha)=\frac{2(p-1)}{p}$.

\noindent \textbf{General case}. Suppose  $l=p_1^{r_1}p_2^{r_2}\ldots p_k^{r_k}\neq 6,12,10,14$, where the $p_i$ are pairwise different primes and each $r_i\ge 1$. For each $1\le i\le k$, we define $w_{p_i^{r_i}}(\alpha)=\frac{\prod\limits_{j\neq i} p_j^{r_j}}{k}$, unless $p_i^{r_i}=2$ and $2\nparallel n$. In that case, $w_2(\alpha)=0$, while $w_4(\alpha)=\frac{\prod\limits_{j\neq i} p_j^{r_j}}{2k}$. In particular, if $l=p^r\ge 3$ is a prime power, $w_{p^r}(\alpha)=1$.

Note that, as we required, the sum $\sum\limits_{p^r|n} w_{p^r}(\alpha)$ over all the prime powers dividing~$n$ is the length $l$ of $\alpha$. Note also that if $l=p_1^{r_1}p_2^{r_2}\ldots p_k^{r_k}$, then $w_{p^r}(\alpha)\neq 0$ only if $p^r=p_i^{r_i}$ for some $1\le i\le k$ or $p^r=4$. 

\noindent\textit{Required core: the prime-power weight theorem and its complete proof.}
\begin{teo} \label{main}
Let $n\ge 1$. Let $P$ be a poset with $\aut (P)$ cyclic of order $n$ generated by $g$. Let $p^r$ be a prime power which exactly divides $n$. If $p^r\neq 2,4$ then \mbox{$\sum\limits_{\alpha \in \sigma} w_{p^r}(\alpha)\ge b(p^r)$}. 
If~$3\nparallel n$ and~$p^r=2$ or~$p^r=4$, $\sum\limits_{\alpha \in \sigma} w_{p^r}(\alpha)\ge b(p^r)$ as well.
If~$3\parallel n$ and~$2\parallel n$, $\sum\limits_{\alpha \in \sigma} (2w_{2}(\alpha)+3w_{3}(\alpha))\ge 2b(2)+3b(3)=11$. Finally, if $3\parallel n$ and~$4\parallel n$, $\sum\limits_{\alpha \in \sigma} (4w_{4}(\alpha)+3w_{3}(\alpha))\ge 4b(4)+3b(3)-1=20$.
\end{teo}

\begin{proof}
If $p^r\neq 2,3,4,5,7$, by Lemma~\ref{dos}, there are at least two cycles of length divisible by $p^r$. By hypothesis their lengths are not multiples of $p^{r+1}$. But if $\alpha$ is a cycle of $\sigma$ whose length is a multiple of $p^r$, then $w_{p^r}(\alpha)\ge 1$. Indeed, the weights in $\alpha$ are assigned according to the General case. If the length of $\alpha$ is $l=p_1^{r_1}p_2^{r_2}\ldots p_k^{r_k}$, we can assume $p^r=p_1^{r_1}$ and then $w_{p^r}(\alpha)=\frac{\prod\limits_{j=2}^k p_j^{r_j}}{k}\ge \frac{2^{k-1}}{k}\ge 1$. Thus, $\sum\limits_{\alpha \in \sigma} w_{p^r}(\alpha)\ge 2= b(p^r)$.

Suppose now $p^r=5$. If $\alpha$ is a cycle of $\sigma$ of length $l=5$, then $w_5(\alpha)=1$. If $l=10$, then $w_5(\alpha)=\frac{8}{5}\ge \frac{3}{2}$ (Exception 10-14). If $l=5s$ with $s=p_2^{r_2}p_3^{r_3}\ldots p_k^{r_k}\ge 3$ not divisible by~$5$, then either $k=2$, or $k\ge 3$. In the first case $w_5(\alpha)=\frac{s}{2}\ge \frac{3}{2}$, and in the second case $w_5(\alpha)=\frac{\prod\limits_{j=2}^k p_j^{r_j}}{k}\ge \frac{2^{k-2}.3}{k}\ge 2\ge \frac{3}{2}$. 

By Lemma~\ref{dos}, there are at least two cycles of length divisible by $5$ (and not by~$5^2$). Suppose first there are exactly two such cycles, $\alpha$ and $\alpha'$. None of them can be of length~$5$ by Lemma~\ref{lema357}. Thus $w_5(\alpha)+w_5(\alpha')\ge 2. \frac{3}{2}= 3=b(5)$. Finally, if there are at least three cycles in $\sigma$ of length divisible by $5$, then $\sum\limits_{\alpha \in \sigma} w_{5}(\alpha)\ge 3= b(5)$.

The case $p^r=7$ is similar to the previous one, with the observation that for length $l=14$, $w_7(\alpha)=\frac{12}{7}\ge \frac{3}{2}$ (Exception 10-14). So, also in this case~$\sum\limits_{\alpha \in \sigma} w_{7}(\alpha)\ge 3= b(7)$.

Let $p^r=3$. If the length of a cycle $\alpha$ in $g$ is $l=3$, $w_3(\alpha)=1$. If $l=6$, $w_3(\alpha)=2$ (Exception 6). If $l=12$, $w_3(\alpha)=4$ (Exception 12). If $l=3s$ with $s=p_2^{r_2}p_3^{r_3}\ldots p_k^{r_k}\ge 5$, then either $k=2$, or $k\ge 3$. In the first case $w_3(\alpha)=\frac{s}{2}\ge \frac{5}{2}$, and in the second case~$w_3(\alpha)=\frac{\prod\limits_{j=2}^k p_j^{r_j}}{k}\ge \frac{2^{k-2}.3}{k}\ge 2$.

By Lemma~\ref{dos} there are at least two cycles in $\sigma$ of length divisible by $3$ (and not by $3^2$). Suppose first there are exactly two such cycles $\alpha$ and $\alpha'$. None of them can have length $3$ by Lemma~\ref{lema357}. Then $w_3(\alpha)+w_3(\alpha')\ge 2.2=4\ge 3=b(3)$. Finally, if there are at least three cycles in $\sigma$ of length divisible by $3$, then $\sum\limits_{\alpha \in \sigma} w_{3}(\alpha)\ge 3= b(3)$. Note that $\sum\limits_{\alpha \in \sigma} w_{3}(\alpha)\ge 4$ unless there are exactly three cycles of length $3$ and no other cycle of length divisible by $3$.

We analyze now the case that $3\nparallel n$ and $p^r=2$ or $4$. In the first situation, there is at least one cycle $\alpha$ of even length $l$ (not divisible by $4$). If $l=2$, $w_2(\alpha)=1$ (General case). If $l=6$, $w_2(\alpha)=3$ (Exception 6). If $l=10$ or $l=14$, then $w_2(\alpha)=1$ (Exception 10-14). If $l=2s$ with $s=p_2^{r_2}p_3^{r_3}\ldots p_k^{r_k}\neq 1,3,5,7$ (odd), then $w_2(\alpha)=\frac{\prod\limits_{j=2}^k p_j^{r_j}}{k}\ge \frac{3^{k-1}}{k}\ge \frac{3}{2}$. Thus $\sum\limits_{\alpha \in \sigma} w_{2}(\alpha)\ge 1= b(2)$. We consider the second situation, $p^r=4$. If $\alpha$ has length $l=4$, then $w_4(\alpha)=1$. If $l=12$, $w_4(\alpha)=3$ (Exception 12). If $l=4s$ with $s=p_2^{r_2}p_3^{r_3}\ldots p_k^{r_k}\ge 5$ (odd), then $k=2$ or $k\ge 3$. For $k=2$ we have $w_4(\alpha)=\frac{s}{2}\ge \frac{5}{2}$. For $k\ge 3$, $w_4(\alpha)\ge \frac{3^{k-1}}{k}\ge 3$. By Lemma~\ref{dos}, $\sigma$ contains at least two cycles of lengths divisible by $4$ (and not by $8$). Suppose first there are exactly two such cycles, $\alpha$ and $\alpha'$, of lengths $l,l'$. If $l=l'=4$, then by Lemma~\ref{lema4}, there exists a third and a fourth cycle $\beta, \beta'$ of lengths $2m$ and $2m'$ for some odd $m,m'$. The weights $w_4(\beta)$ that we obtain for each $m$ are the halves of the weights that we obtained for~$2$ in cycles of the same length when $2\parallel n$. Namely, if $m=1$, $w_4(\beta)=\frac{1}{2}$ (General case); if $m=3$, $w_4(\beta)=\frac{3}{2}$ (Exception 6); if $m=5,7$, $w_4(\beta)=\frac{1}{2}$ (Exception 10-14); if $m=p_2^{r_2}p_3^{r_3}\ldots p_k^{r_k}\neq 1,3,5,7$ then $w_4(\beta)\ge \frac{3^{k-1}}{2k}\ge \frac{3}{4}$ (General case).

The same happens with $\beta'$. Thus $w_4(\alpha)+w_4(\alpha')+w_4(\beta)+w_4(\beta')\ge 1+1+\frac{1}{2}+\frac{1}{2}=3=b(4)$. If instead $l=4$ and $l'=12$, then $w_4(\alpha)+w_4(\alpha')=1+3=4> 3$. If $l=4$ and $l'=4s$ for some odd $s\ge 5$, then $w_4(\alpha)+w_4(\alpha')\ge 1+\frac{5}{2}>3$. If both $l$ and $l'$ are greater than $4$, then $w_4(\alpha)+w_4(\alpha')\ge \frac{5}{2}+\frac{5}{2}>3$. Finally, if there are at least three cycles of length divisible by $4$, then $\sum\limits_{\alpha \in \sigma} w_4(\alpha)\ge 3$. Thus, in any case $\sum\limits_{\alpha \in \sigma} w_4(\alpha)\ge 3=b(4)$.

It only remains to analyze the case $3\parallel n$ and $2\parallel n$ and the case $3\parallel n$ and $4\parallel n$.        
If $3\parallel n$ and $2\parallel n$, recall that we have already proved that $\sum\limits_{\alpha \in \sigma} w_{3}(\alpha)\ge 4$ or there are exactly three cycles of length $3$ and no other cycle of length divisible by $3$. In the first case $\sum\limits_{\alpha \in \sigma} (2w_{2}(\alpha)+3w_{3}(\alpha))\ge \sum\limits_{\alpha \in \sigma} 3w_{3}(\alpha)\ge 12$. In the second case, there exists a cycle $\beta$ in $\sigma$ of even length $m\neq 6$, so $w_2(\beta)\ge 1$. Thus $\sum\limits_{\alpha \in \sigma} (2w_{2}(\alpha)+3w_{3}(\alpha))\ge 2.1+3.3= 11$. 

The last case is $3\parallel n$ and $4\parallel n$.
Note that if there are no cycles of length $6$ nor~$12$ in $\sigma$, then the computation $\sum\limits_{\alpha \in \sigma} w_4(\alpha)\ge 3$ remains valid as Exceptions 6 and 12 do not occur. Thus $\sum\limits_{\alpha \in \sigma} (3w_{3}(\alpha)+4w_{4}(\alpha))\ge 3.3+4.3=21>20$. If there are at least two $12$-cycles, then $\sum\limits_{\alpha \in \sigma} (3w_{3}(\alpha)+4w_{4}(\alpha))\ge 2.3.4=24>20$. If there is no $12$-cycle in $\sigma$ and $\sum\limits_{\alpha \in \sigma} w_4(\alpha)< 3$, then we must be in the case that there is a $6$-cycle. This already implies $\sum\limits_{\alpha \in \sigma} w_3(\alpha)\ge 4$, while the existence of two cycles of length divisible by~$4$ implies $\sum\limits_{\alpha \in \sigma} w_4(\alpha)\ge 2$. Thus $\sum\limits_{\alpha \in \sigma} (3w_{3}(\alpha)+4w_{4}(\alpha))\ge 3.4+4.2=20$.

Thus we may assume $\sigma$ has a unique $12$-cycle. By Lemma~\ref{dos} there is another cycle of length divisible by $4$, so $\sum\limits_{\alpha \in \sigma} w_4(\alpha)\ge 1$. On the other hand, $\sum\limits_{\alpha \in \sigma} w_{3}(\alpha)\ge 4+2=6$, as the weight of $3$ in a $12$-cycle is $4$ and by Lemmas~\ref{dos} and~\ref{lema357} there are either two more cycles of lengths divisible by $3$ or just one, but of length not $3$. Thus $\sum\limits_{\alpha \in \sigma} (3w_{3}(\alpha)+4w_{4}(\alpha))\ge 3.6+4.1=22>20$.
\end{proof}

\section*{Counted claim: original Corollary 4.2 with the complete author proof}
\begin{coro} \label{maincoro}
Let $n=p_1^{r_1}p_2^{r_2}\ldots p_k^{r_k}$, where the $p_i$ are pairwise different primes and~$r_i\ge 1$ for every $i$. Then the minimum number $\beta(\Z_n)$ of points in a poset with cyclic automorphism group of order $n$ is $\sum\limits_{i=1}^k b(p_i^{r_i})p_i^{r_i}-1$ if $3\parallel n$ and $4\parallel n$, and $\sum\limits_{i=1}^k b(p_i^{r_i})p_i^{r_i}$ otherwise.
\end{coro}

\begin{proof}
If $P$ is a poset with $\aut (P)\simeq \Z_n$ generated by $g$, then the number of points in $P$ is at least $\sum\limits_{\alpha \in g} |\alpha|=$ $\sum\limits_{\alpha \in g} \sum\limits_{p^r|n}w_{p^r}(\alpha)p^r\ge \sum\limits_{i=1}^k (\sum\limits_{\alpha \in g} w_{p_i^{r_i}}(\alpha))p_i^{r_i}$. If both $3$ and $4$ exactly divide $n$, by Theorem~\ref{main} this is $\sum\limits_{p_i^{r_i}\neq 3,4} (\sum\limits_{\alpha \in g} w_{p_i^{r_i}}(\alpha))p_i^{r_i}+\sum\limits_{\alpha \in g} (3w_{3}(\alpha)+4w_4(\alpha))\ge \sum\limits_{p_i^{r_i}\neq 3,4} b(p_i^{r_i})p_i^{r_i}+3b(3)+4b(4)-1=\sum\limits_{i=1}^k b(p_i^{r_i})p_i^{r_i}-1$. Otherwise, the bound is one more than this number. The bound is attained by Theorem~\ref{teoejemplos}.
\end{proof}

\noindent\textit{Excluded material (recorded, not claimed).} The introduction's history of the analogous graph problem, its unnumbered duplicate announcement of the same corollary and its description of the paper structure are not reproduced, and no image, table or other asset of the original is redistributed. The two external results invoked in the excluded history, the Birkhoff and Babai poset constructions and the Frucht construction quoted in the proof of Proposition 2.3, are recorded only through their exact published citations.
\begin{thebibliography}{99}
\bibitem[4]{Bab2} L. Babai, ``Finite digraphs with given regular automorphism groups'', \emph{Period. Math. Hungar.} \textbf{11} (1980), no. 4, 257--270, https://doi.org/10.1007/BF02107568.
\bibitem[6]{Bir} G. Birkhoff, ``On groups of automorphisms'', \emph{Rev. Un. Mat. Argentina} \textbf{11} (1946), 155--157.
\bibitem[9]{Fru50} R. Frucht, ``On the construction of partially ordered systems with a given group of automorphisms'', \emph{Amer. J. Math.} \textbf{72} (1950), 195--199, https://doi.org/10.2307/2372146.
\end{thebibliography}
\end{document}
